\chapter{Covering, maximal estimates, and Lebesgue points}\label{sec:measure-differentiation}
Small-ball averages recover a function almost everywhere, but can behave very differently on the support of a singular measure. We prove the Euclidean Lebesgue-measure statements directly from a covering argument. The source is \cite{Sim18}, Chapter 1, Lemma 3.4 and Sections 3--4, printed pp.~15--24. Differentiation relative to an arbitrary Radon measure requires additional covering theory and is not asserted here. Balls in this section are open.

\begin{lemma}[Fivefold covering for a bounded family of Euclidean balls]\label{lem:bounded-fivefold-covering}
Let $\mathcal B$ be a family of balls of positive radii in $\mathbb R^n$, with bounded centers and uniformly bounded radii. There is a finite or countable pairwise disjoint subfamily $B_j=B_{r_j}(x_j)$ such that
\[
 \bigcup_{B\in\mathcal B}B\subset\bigcup_j B_{5r_j}(x_j).
\]
Moreover every original ball intersects a selected ball whose radius is at least half its own radius. This is a bounded Euclidean version of \cite{Sim18}, Lemma 3.4.
\end{lemma}
\begin{proof}
If no balls remain, stop. Otherwise select a remaining ball whose radius is greater than half the supremum of the remaining radii, and discard all remaining balls intersecting it. Continue. The selected balls are disjoint and lie in a fixed bounded region. If the process is infinite, their volumes have finite sum, so their radii tend to zero. An original ball of radius $r>0$ cannot remain forever: while it remains, every selected radius is greater than $r/2$, contradicting that limit. Thus it is selected or discarded at some step. If it intersects $B_{r_j}(x_j)$ at that step, its radius is at most $2r_j$, so each of its points lies within distance $r_j+2r\leq5r_j$ of $x_j$. This proves both assertions, with the open-ball inequalities giving containment in the open dilation.
\end{proof}

\begin{definition}[Centered maximal functions]\label{def:centered-maximal-function}
For a nonnegative Radon measure $\nu$ on $\mathbb R^n$ define
\[
 M\nu(x)=\sup_{r>0}\frac{\nu(B_r(x))}{\omega_n r^n}.
\]
For $f\in L^1_{\rm loc}(\mathbb R^n)$ write $Mf=M(|f|\,dx)$. The value may be infinite. These are measurable functions: for fixed $r$, Fatou's lemma makes $x\mapsto\nu(B_r(x))$ lower semicontinuous, and the supremum retains this property. This maximal quantity controls the upper density defined in \cite{Sim18}, Chapter 1, equation (3.1).
\end{definition}

\begin{proposition}[The weak maximal estimate for finite measures]\label{prop:measure-maximal-weak}
If $\nu$ is a finite nonnegative Radon measure on $\mathbb R^n$ and $\lambda>0$, then
\[
 \mathcal L^n(\{M\nu>\lambda\})\leq\frac{5^n}{\lambda}\nu(\mathbb R^n).
\]
In particular $M\nu<\infty$ almost everywhere for Lebesgue measure. The constant is not claimed to be sharp.
\end{proposition}
\begin{proof}
For each $x\in\{M\nu>\lambda\}\cap B_R(0)$ choose a ball centered at $x$ with $\nu(B_r(x))>\lambda\omega_n r^n$. Its radius is bounded in terms of $\nu(\mathbb R^n)$ and $\lambda$. Lemma \ref{lem:bounded-fivefold-covering} gives disjoint selected balls covering the set after fivefold dilation. Therefore its Lebesgue measure is at most
\[
 5^n\sum_j\omega_n r_j^n
 \leq\frac{5^n}{\lambda}\sum_j\nu(B_{r_j}(x_j))
 \leq\frac{5^n}{\lambda}\nu(\mathbb R^n).
\]
Let $R\to\infty$. Letting $\lambda\to\infty$ then shows that the infinite-value set is null. If the measure is zero, the assertions are immediate. This is the maximal-estimate application of the covering principle in \cite{Sim18}, Lemma 3.4.
\end{proof}

\begin{theorem}[A strong maximal bound above the endpoint]\label{thm:maximal-strong-lp}
For $1<p<\infty$ and $f\in L^p(\mathbb R^n)$,
\[
 \|Mf\|_p^p\leq\frac{2^p5^n p}{p-1}\|f\|_p^p.
\]
This consequence of the covering estimate associated with \cite{Sim18}, Lemma 3.4, is proved here by integrating distribution functions.
\end{theorem}
\begin{proof}
Fix $\lambda>0$. Split $|f|$ into the part bounded by $\lambda/2$ and the tail $g_\lambda=|f|\mathbf1_{\{|f|>\lambda/2\}}$. The first part has maximal function at most $\lambda/2$. The tail is integrable since $|f|\leq(\lambda/2)^{1-p}|f|^p$ on its support. Subadditivity and Proposition \ref{prop:measure-maximal-weak} give
\[
 \mathcal L^n(\{Mf>\lambda\})
 \leq\frac{2\cdot5^n}{\lambda}\int_{\{|f|>\lambda/2\}}|f|\,dx.
\]
The identity $a^p=\int_0^a p\lambda^{p-1}\,d\lambda$, followed by Tonelli's theorem, gives
\[
 \int(Mf)^p\leq2\cdot5^n p\int |f(x)|
       \left(\int_0^{2|f(x)|}\lambda^{p-2}\,d\lambda\right)dx.
\]
Evaluation of the inner integral proves the claim. Tonelli applies before finiteness is known because the integrands are nonnegative. Local integrability of $f$, needed to define $Mf$, follows from Holder's inequality on bounded sets.
\end{proof}

\begin{definition}[Lebesgue points]\label{def:lebesgue-point}
For a finite-valued representative of $f\in L^1_{\rm loc}(U)$, a point $x\in U$ is a Lebesgue point if
\[
 \lim_{r\downarrow0}\frac{1}{\omega_n r^n}
       \int_{B_r(x)}|f(y)-f(x)|\,dy=0,
\]
with radii small enough that the ball lies in $U$. Altering values on a null set can change individual Lebesgue points but not the almost-everywhere conclusion below. This is the absolute-oscillation form of differentiation associated with \cite{Sim18}, Corollary 4.10.
\end{definition}

\begin{theorem}[Lebesgue differentiation with absolute oscillation]\label{thm:lebesgue-differentiation}
If $f\in L^1_{\rm loc}(U)$, almost every point of $U$ is a Lebesgue point of any representative finite almost everywhere. Consequently
\[
 \lim_{r\downarrow0}\frac{1}{\omega_n r^n}\int_{B_r(x)}f(y)\,dy=f(x)
 \quad\text{for almost every }x\in U.
\]
Compare \cite{Sim18}, Corollary 4.10. Here the proof is restricted to Lebesgue measure on a Euclidean open set.
\end{theorem}
\begin{proof}
First take $U=\mathbb R^n$ and $f\in L^1$. Denote the limsup in Definition \ref{def:lebesgue-point} by $Df(x)$. For $g\in C_c(\mathbb R^n)$, continuity gives $Dg=0$, so the triangle inequality implies
\[
 Df(x)\leq M(f-g)(x)+|f(x)-g(x)|.
\]
The weak maximal estimate and Chebyshev's inequality yield, for $a>0$,
\[
 \mathcal L^n(\{Df>2a\})\leq\frac{5^n+1}{a}\|f-g\|_1.
\]
These sets are measurable: for each point with finite $f(x)$ the ball integrals vary continuously with the positive radius, so the limsup can be taken over rational radii; joint measurability follows from integration of the measurable integrand. By Theorem \ref{thm:radon-lp-density}, choose $g$ with arbitrarily small $L^1$ error. Thus $\{Df>2a\}$ is null for each $a>0$. A countable sequence $a\downarrow0$ proves $Df=0$ almost everywhere.

For the local assertion, choose compact sets $K_m$ with $K_m\subset\operatorname{int}K_{m+1}$ exhausting $U$. Multiply $f$ by a compactly supported cutoff equal to one near $K_m$, and extend by zero outside $U$. This gives an $L^1(\mathbb R^n)$ function agreeing with $f$ near $K_m$. The global result applies on each $K_m$; their countable union gives the local assertion. Finally the absolute value of the difference between the average of $f$ and $f(x)$ is bounded by the average absolute oscillation.
\end{proof}

\begin{corollary}[The Lebesgue density theorem for measurable sets]\label{cor:lebesgue-density-theorem}
For every Lebesgue measurable $E\subset\mathbb R^n$,
\[
 \lim_{r\downarrow0}\frac{\mathcal L^n(E\cap B_r(x))}{\omega_n r^n}
 =\mathbf1_E(x)\quad\text{for almost every }x\in\mathbb R^n.
\]
Neither finite total measure nor a regular boundary is assumed. See \cite{Sim18}, Corollary 3.10.
\end{corollary}
\begin{proof}
The indicator belongs to $L^1_{\rm loc}$. Apply Theorem \ref{thm:lebesgue-differentiation} to it. At a point in $E$, its absolute oscillation average is the complementary volume fraction; outside $E$, it is the volume fraction of $E$ itself. The conclusion is almost everywhere in ambient Lebesgue measure, and makes no assertion at every topological boundary point.
\end{proof}

\begin{proposition}[Approximate continuity at a Lebesgue point]\label{prop:lebesgue-approximate-continuity}
If $x$ is a Lebesgue point of $f$ and $a>0$, then
\[
 \lim_{r\downarrow0}
 \frac{\mathcal L^n(\{y\in B_r(x):|f(y)-f(x)|>a\})}{\omega_n r^n}=0.
\]
\end{proposition}
\begin{proof}
Chebyshev's inequality bounds the displayed fraction by $a^{-1}$ times the average of $|f(y)-f(x)|$ over the ball. That average tends to zero. Thus values far from $f(x)$ occupy a vanishing fraction of small balls; ordinary pointwise continuity is a stronger condition.
\end{proof}

\begin{proposition}[Mollifiers converge at every Lebesgue point]\label{prop:mollifiers-at-lebesgue-points}
Let $f\in L^1_{\rm loc}(\mathbb R^n)$ and let $\rho$ be a bounded measurable function supported in $B_1$, with $\int\rho=1$. Put $\rho_\varepsilon(z)=\varepsilon^{-n}\rho(z/\varepsilon)$. At every Lebesgue point $x$ of $f$,
\[
 (\rho_\varepsilon*f)(x)\longrightarrow f(x).
\]
Nonnegativity and smoothness of $\rho$ are not needed for this pointwise assertion.
\end{proposition}
\begin{proof}
Subtract $f(x)\int\rho_\varepsilon=f(x)$ from the convolution and estimate its absolute value by
\[
 \|\rho\|_\infty\varepsilon^{-n}
       \int_{B_\varepsilon(x)}|f(y)-f(x)|\,dy.
\]
This is $\omega_n\|\rho\|_\infty$ times a quantity tending to zero. The smooth-density approximation of measures in Proposition \ref{prop:radon-mollification} concerns compact-test convergence; the present assertion describes values of mollified functions at individual Lebesgue points.
\end{proof}

\begin{proposition}[A singular measure has zero ambient density almost everywhere]\label{prop:singular-ambient-density-zero}
Let $\nu$ be a nonnegative Radon measure on $\mathbb R^n$ singular with respect to $\mathcal L^n$. Then
\[
 \lim_{r\downarrow0}\frac{\nu(B_r(x))}{\omega_n r^n}=0
 \quad\text{for }\mathcal L^n\text{-almost every }x.
\]
This is a special case of the upper-density conclusions in \cite{Sim18}, Theorem 3.8 and Remark 3.9.
\end{proposition}
\begin{proof}
First assume $\nu$ is finite, and choose a Borel null set $N$ carrying all its mass. Given $\varepsilon>0$, inner regularity gives compact $K\subset N$ with $\nu(\mathbb R^n\setminus K)<\varepsilon$. For $x\notin K$, all sufficiently small balls avoid $K$, so the upper density of $\nu$ at $x$ is at most $M(\nu\llcorner(\mathbb R^n\setminus K))(x)$. Proposition \ref{prop:measure-maximal-weak} implies that the set where this upper density exceeds $a>0$ has Lebesgue measure at most $5^n\varepsilon/a$, since $K$ is null. Let $\varepsilon\downarrow0$ and then take countably many $a\downarrow0$.

For locally finite $\nu$, apply the finite result to its restriction to $B_{m+1}(0)$ and consider points of $B_m(0)$. Restriction does not change sufficiently small balls there. Taking the union over positive integers $m$ proves the assertion. Restrictions to these Borel sets remain Radon measures.
\end{proof}

\begin{corollary}[Ball densities identify an absolutely continuous part]\label{cor:ball-density-ac-part}
Suppose a nonnegative Radon measure is written $\mu=f\,dx+\nu$ with $0\leq f\in L^1_{\rm loc}(\mathbb R^n)$ and $\nu\perp\mathcal L^n$. Then
\[
 \lim_{r\downarrow0}\frac{\mu(B_r(x))}{\omega_n r^n}=f(x)
 \quad\text{for }\mathcal L^n\text{-almost every }x.
\]
\end{corollary}
\begin{proof}
Add the limits supplied by Theorem \ref{thm:lebesgue-differentiation} and Proposition \ref{prop:singular-ambient-density-zero}, outside the union of their null exceptional sets. The decomposition is a hypothesis of this statement; existence of the Lebesgue decomposition and the Radon--Nikodym theorem are not proved by this argument. Compare \cite{Sim18}, Section 4, especially Theorem 4.18.
\end{proof}

\begin{example}[The maximal operator need not preserve integrability]\label{ex:maximal-endpoint-failure}
Let $f=\mathbf1_{B_1(0)}$ on $\mathbb R^n$. Then $f\in L^1$ but $Mf\notin L^1$.
\end{example}
\begin{proof}
For $|x|\geq2$, the ball $B_{|x|+1}(x)$ contains $B_1(0)$, so
\[
 Mf(x)\geq (|x|+1)^{-n}\geq(3|x|/2)^{-n}.
\]
The integral of $|x|^{-n}$ over $\{|x|\geq2\}$ diverges: polar coordinates reduce it to a positive constant times $\int_2^\infty r^{-1}\,dr$. This separates the weak estimate in Proposition \ref{prop:measure-maximal-weak} from the strong estimates for $p>1$ in Theorem \ref{thm:maximal-strong-lp}.
\end{proof}
